% !TEX root = Lecture_Koszul.tex
\section{Graded algebras, modules, and duals}

\subsection{Algebras and modules}

\begin{definition}
A \defn{$\Bbbk$-algebra} is a $\Bbbk$-vector space with an associative unital multiplication that is \defn{bilinear}, meaning linear in each variable.
\end{definition}

A right $A$-module is understood to be a $\Bbbk$-vector space $M$ with a bilinear action satisfying $(ma)b=m(ab)$ and $m1_A=m$ for $m\in M$ and $a,b\in A$.
A module homomorphism $f:M\to N$ is linear and satisfies $f(ma)=f(m)a$.
A homomorphism of $\Bbbk$-algebras is a linear map preserving multiplication and the identity.

\begin{example}\label{ex:triangular-algebra}
(1) The matrix ring $\Mat_n(\Bbbk)$ is a $\Bbbk$-algebra.
The space of $1\times n$ row vectors is a right $\Mat_n(\Bbbk)$-module by matrix multiplication.

(2) Let $\Bbbk A_n$ be the \defn{upper triangular matrix ring}, which is also a $\Bbbk$-algebra.
\[
\Bbbk A_n = \begin{pmatrix}
\Bbbk & \Bbbk & \cdots & \Bbbk \\
0 & \Bbbk & \cdots & \Bbbk \\
0 & 0 & \ddots & \vdots \\
0 & \cdots & 0 & \Bbbk
\end{pmatrix}
\]
For those familiar with path algebras, $\Bbbk A_n$ is the path algebra of a \defn{linearly oriented $A_n$ quiver}, which explains the notation. We use matrix calculations, so no knowledge of quivers is required.
Write $e_{ij}$ for the matrix with $1$ in entry $(i,j)$ and zero elsewhere.
Our convention composes paths from left to right: $e_{ij}$ represents $i\to j$, and $e_{ij}e_{j\ell}=e_{i\ell}$.
\end{example}

\subsection{Categories and functors}\label{subsec:category-language}

We organise modules together with their homomorphisms. A category records both: specifying the objects alone does not specify a category.

\begin{definition}\label{def:category}
A \defn{category} $\CC$ consists of objects, a set $\Hom_{\CC}(X,Y)$ of \defn{morphisms} $f:X\to Y$ for each pair of objects, and composition sending $f:X\to Y$ and $g:Y\to Z$ to $g\circ f:X\to Z$.
Each object $X$ has an identity $\id_X:X\to X$, and
\[
h\circ(g\circ f)=(h\circ g)\circ f,
\qquad \id_Y\circ f=f=f\circ\id_X
\quad\text{for all composable }f,g,h.
\]
\end{definition}

A morphism $f:X\to Y$ is an \defn{isomorphism} if some $g:Y\to X$ satisfies $g\circ f=\id_X$ and $f\circ g=\id_Y$; we then write $X\cong Y$.
A \defn{full subcategory} of $\CC$ consists of chosen objects of $\CC$, with all morphisms between them and the same composition.

\begin{example}\label{ex:module-categories}
For a $\Bbbk$-algebra $A$, we use the notation
\[
\begin{aligned}
\Mod A&:=\text{category of right $A$-modules and $A$-linear maps},\\
\mod A&:=\text{full subcategory of finitely generated right $A$-modules}.
\end{aligned}
\]
In either category, $\Hom_A(M,N)$ is the space of $A$-linear maps $M\to N$; composition is the usual composition of maps.
For $A=\Bbbk$, these are the categories of all and of finite-dimensional vector spaces, respectively.

Note that, if $\dim_\Bbbk A<\infty$, then $M\in \mod A \;\Leftrightarrow\; M$ is finite-dimensional ($\because$ finite generation implies that $M$ is a quotient of $A^{\oplus n}$ some $n$, and a finite $\Bbbk$-basis of $M$ is clearly a finite generating set.)
\end{example}

\begin{example}\label{ex:add-category}
For $M\in\mod A$, let $\add M:=\add_{\mod A}(M)$ be the full subcategory with
\[
N\in\add M
\quad\Longleftrightarrow\quad
N\oplus N'\cong M^{\oplus r}
\quad\text{for some }N'\in\mod A\text{ and integer }r\geq0.
\]
Thus $\add M$ consists of direct summands of finite direct sums of $M$.
\end{example}

In modern representation theory, remembering $A$ and $\add A$ holds the same categorical information is crucial in utilising cateogircal technique in representation theory; this allows one to push many results to a more abstract framework.

\begin{definition}\label{def:functor}
A \defn{functor} $F:\CC\to\CD$ assigns an object $F(X)$ to each object $X$ and a morphism $F(f):F(X)\to F(Y)$ to each $f:X\to Y$, such that
\[
F(\id_X)=\id_{F(X)},\qquad F(g\circ f)=F(g)\circ F(f)
\quad\text{for all composable }f,g.
\]
\end{definition}

For example, forgetting the $A$-action gives a functor $\Mod A\to\Mod\Bbbk$: it keeps the underlying vector spaces and linear maps.
The \defn{identity functor} $\id_{\CC}$ leaves every object and morphism unchanged; functors compose by applying one after the other.
Every functor preserves isomorphisms, since $F(f^{-1})$ is inverse to $F(f)$.

\phantomsection\label{not:opposite-category}
The \defn{opposite category} $\CC^\op$ has the same objects as $\CC$, with $\Hom_{\CC^\op}(X,Y)=\Hom_{\CC}(Y,X)$ and composition reversed.
A \defn{contravariant functor} from $\CC$ to $\CD$ means a functor $\CC^\op\to\CD$.
Thus it sends $f:X\to Y$ in $\CC$ to $F(f):F(Y)\to F(X)$ and satisfies $F(g\circ f)=F(f)\circ F(g)$.
Duals will provide our main examples.

\begin{definition}\label{def:natural-transformation}
For functors $F,G:\CC\to\CD$, a \defn{natural transformation} $\eta:F\to G$ is a family of morphisms $\eta_X:F(X)\to G(X)$ such that
\begin{equation}\label{eq:naturality-square}
\xymatrix@C=45pt@R=20pt{
F(X) \ar[r]^{\eta_X} \ar[d]_{F(f)} & G(X) \ar[d]^{G(f)}\\
F(Y) \ar[r]_{\eta_Y} & G(Y)
}
\qquad
G(f)\circ\eta_X=\eta_Y\circ F(f)
\quad\text{for every }f:X\to Y.
\end{equation}
If every $\eta_X$ is an isomorphism, $\eta$ is a \defn{natural isomorphism}, written $F\cong G$.
\end{definition}

The square in \eqref{eq:naturality-square} is said to \defn{commute}: its two routes from $F(X)$ to $G(Y)$ give the same morphism.
We often say a family of morphisms (resp. isomorphisms) in a category $\CC$ is \defn{natural in $M\in\CC$} to mean that it defines a natural transformation (resp. isomorphism).

For example, for fixed $c\in\Bbbk$, the maps $c\id_M:M\to M$ define a natural transformation $\id_{\Mod A}\to\id_{\Mod A}$ because $f(cm)=cf(m)$ for every module homomorphism $f:M\to N$.
It is a natural isomorphism when $c\neq0$.

A functor $F:\CC\to\CD$ is an \defn{equivalence} if there is a functor $G:\CD\to\CC$ with natural isomorphisms $GF\cong\id_{\CC}$ and $FG\cong\id_{\CD}$.
A \defn{duality} between $\CC$ and $\CD$ is an equivalence $\CC^\op\to\CD$.
Thus an equivalence recovers objects and maps up to natural isomorphism; a duality also reverses the direction of maps.

\begin{definition}\label{def:exact-sequence}
Let $A$ be a $\Bbbk$-algebra, and let $L,M,N\in\Mod A$, with $i:L\to M$ and $p:M\to N$ module homomorphisms.
The sequence $L\xrightarrow{i}M\xrightarrow{p}N$ is \defn{exact at $M$} if $\Im i=\Ker p$.  
A sequence of module homomorphisms is \defn{exact} if it is exact at every term with an incoming and an outgoing map.
A \defn{short exact sequence} is an exact sequence
\begin{equation}\label{eq:short-exact-sequence}
0\longrightarrow L\xrightarrow{i}M\xrightarrow{p}N\longrightarrow0,
\qquad
\Ker i=0,\quad\Im i=\Ker p,\quad\Im p=N.
\end{equation}
\end{definition}


\begin{definition}\label{def:exact-functor}\label{def:exact-contravariant}
Let $A,B$ be $\Bbbk$-algebras.
A functor $F:\Mod A\to\Mod B$ is \defn{exact} if $F(0)=0$ and every short exact sequence \eqref{eq:short-exact-sequence} gives a short exact sequence
\[
0\longrightarrow F(L)\xrightarrow{F(i)}F(M)\xrightarrow{F(p)}F(N)\longrightarrow0.
\]
\end{definition}

An exact functor preserves injective and surjective maps: apply it to
$0\to L\xrightarrow{i}M\to M/i(L)\to0$ and $0\to\Ker p\to M\xrightarrow{p}N\to0$, respectively.
For a functor $F:(\Mod A)^\op\to\Mod B$, exactness instead means that \eqref{eq:short-exact-sequence} gives
\[
0\longrightarrow F(N)\xrightarrow{F(p)}F(M)\xrightarrow{F(i)}F(L)\longrightarrow0.
\]
Thus an exact contravariant functor exchanges injective and surjective maps.

\begin{example}[$\Bbbk$-linear duality]\label{ex:linear-duality}
For a vector space $V$, define its \defn{$\Bbbk$-linear dual} by $D(V):=\Hom_\Bbbk(V,\Bbbk)$.
For a linear map $u:V\to W$, set
\[
D(u):D(W)\longrightarrow D(V),\qquad \lambda\longmapsto\lambda\circ u.
\]
The identities $D(\id_V)=\id_{D(V)}$ and $D(v\circ u)=D(u)\circ D(v)$ make $D:(\Mod\Bbbk)^\op\to\Mod\Bbbk$ a functor.
The functor $D$ is exact.
\begin{supplementary}[nobreak=true]
For a short exact sequence \eqref{eq:short-exact-sequence} of vector spaces, $D(p)$ is injective because $p$ is surjective, and $D(i)$ is surjective because a functional on a subspace extends to the whole vector space.
If $\lambda\in D(M)$ satisfies $\lambda\circ i=0$, then
\[
\mu:N\longrightarrow\Bbbk,\qquad p(m)\longmapsto\lambda(m)
\]
is well-defined since $\Ker p=i(L)$, and $\lambda=\mu\circ p=D(p)(\mu)$.
Conversely, $p\circ i=0$ gives $D(i)\circ D(p)=0$, so $\Ker D(i)=\Im D(p)$.
Thus $D$ exchanges injective and surjective linear maps.
\end{supplementary}

For every vector space $V$, evaluation defines a natural injective map
\begin{equation}\label{eq:vector-bidual-evaluation}
\mathrm{ev}_V:V\longrightarrow D(D(V)),\qquad
v\longmapsto(\lambda\mapsto\lambda(v)).
\end{equation}
If $V$ is finite-dimensional, choose a basis $v_1,\ldots,v_n$ and its \defn{dual basis} $v_1^*,\ldots,v_n^*$, defined by $v_i^*(v_j)=\delta_{ij}$.
Then
\[
\mathrm{ev}_V^{-1}:D(D(V))\longrightarrow V,\qquad
\Phi\longmapsto\sum_{i=1}^n\Phi(v_i^*)v_i.
\]
Consequently, $D:(\mod\Bbbk)^\op\to\mod\Bbbk$ is a duality.
\begin{supplementary}
If $v\neq0$, choose $\lambda\in D(V)$ with $\lambda(v)\neq0$; hence $\mathrm{ev}_V$ is injective.
For $u:V\to W$, $v\in V$, and $\lambda\in D(W)$,
\[
D(D(u))(\mathrm{ev}_V(v))(\lambda)=\lambda(u(v))
=\mathrm{ev}_W(u(v))(\lambda),
\]
which proves naturality.
\begin{sloppypar}
In finite dimension, writing $\lambda=\sum_{i=1}^n\lambda(v_i)v_i^*$ gives
$\Phi(\lambda)=\lambda(\sum_{i=1}^n\Phi(v_i^*)v_i)$.
Thus the displayed formula is the inverse of \eqref{eq:vector-bidual-evaluation}, and is independent of the chosen basis.
\end{sloppypar}
\end{supplementary}
\end{example}

\subsection{Semisimple algebras}

\begin{definition}\label{def:semisimple}
A module $L$ is \defn{simple} if $L\neq 0$ and its only submodules are $0$ and $L$.
A module is \defn{semisimple} if it is a direct sum of simple modules.
An algebra $B$ is \defn{semisimple} if it is semisimple as a right module $B_B$ over itself.
The \defn{Jacobson radical} $\rad B$ is the intersection of all maximal right ideals of $B$.
For a right $B$-module $M$, the \defn{radical} $\rad M$ is the intersection of all maximal submodules of $M$, with $\rad M=M$ if there are none.
\end{definition}

We use the following standard ring-theoretic facts without proof.
For a finite-dimensional $\Bbbk$-algebra $B$ and any right $B$-module $M$,
\[
\rad M=M\rad B.
\]
Recall that a \defn{division $\Bbbk$-algebra} is a non-zero $\Bbbk$-algebra whose non-zero elements are invertible.

\begin{theorem}[Artin--Wedderburn]\label{thm:artin-wedderburn}
For a non-zero finite-dimensional $\Bbbk$-algebra $B$,
\[
B\text{ is semisimple}
\quad\Longleftrightarrow\quad
\rad B=0
\quad\Longleftrightarrow\quad
B\cong\prod_{i=1}^n\Mat_{n_i}(D_i),
\]
where $n,n_i\geq 1$ and each $D_i$ is a finite-dimensional division $\Bbbk$-algebra.
If $\Bbbk$ is algebraically closed, then $D_i=\Bbbk$ for every $i$.
For semisimple $B$, every left or right $B$-module is semisimple, and every submodule has a complement. In particular, the definition is independent of the chosen side.
\end{theorem}

\begin{advanced}
For proofs of $\rad M=M\rad B$ and of Theorem~\ref{thm:artin-wedderburn}, see for example \cite{DK94}; \cite{Bez23} gives a short account of semisimple modules and Wedderburn's theorem.
\end{advanced}

\begin{corollary}\label{cor:semisimple-summands}
Every finite-dimensional left or right module over a finite-diemnsional semisimple algebra $S$ is a direct summand of $S^{\oplus n}$ for some $n$, i.e. $M\in \add(S)$ for all $M\in \mod S$.
\end{corollary}

\begin{proof}
Choose a surjection $p:S^{\oplus n}\onto M$.
By Theorem~\ref{thm:artin-wedderburn}, there is a submodule $L$ such that $S^{\oplus n}=\Ker p\oplus L$.
The restriction $p|_L:L\to M$ has kernel $L\cap\Ker p=0$ and image $p(L)=p(S^{\oplus n})=M$.
Hence $L\cong M$, so $S^{\oplus n}\cong\Ker p\oplus M$.
\end{proof}

\begin{example}\label{ex:semisimple-algebras}
The algebras $\Bbbk^{\oplus n}$ and $\Mat_n(\Bbbk)$ are semisimple by Theorem~\ref{thm:artin-wedderburn}.
Upper triangular matrices need not form a semisimple algebra: for $A=\Bbbk A_2$,
\[
\rad A=\Bbbk e_{12}\neq0,\qquad
A/\rad A\cong\Bbbk^{\oplus2}.
\]
\begin{supplementary}
Put $J=\Bbbk e_{12}$, so $J^2=0$.
For a simple right module $L$, simplicity gives $LJ=0$ or $LJ=L$; the second alternative implies $L=LJ^2=0$, a contradiction.
Hence $J$ is contained in every maximal right ideal.
The quotient $A/J\cong\Bbbk^{\oplus2}$ has exactly two maximal right ideals, whose inverse images have intersection
\[
(\Bbbk e_{11}\oplus\Bbbk e_{12})\cap
(\Bbbk e_{12}\oplus\Bbbk e_{22})=J.
\]
\end{supplementary}
\end{example}

\begin{definition}\label{def:idempotents}
An element $e\in A$ is \defn{idempotent} if $e^2=e$.
Idempotents $e,f$ are \defn{orthogonal} if $ef=fe=0$.
A non-zero idempotent is \defn{primitive} if it is not the sum of two non-zero orthogonal idempotents.
A finite family $e_1,\ldots,e_r$ of pairwise orthogonal idempotents is \defn{complete} if $\sum_{i=1}^r e_i=1_A$.
\end{definition}

\begin{remark}\label{rem:idempotent-units}
For orthogonal idempotents $e,f\in A$,
\[
(e+f)A=eA\oplus fA.
\]
In particular, we always have $A\cong eA\oplus (1-e)A$.  
This strengthen the viewpoint of idempotent-as-local-identity: left-multiplying $e$ defines a canonical projection onto the direct summand $eA$ of $A$, and yields a commutative diagram
\[
\xymatrix@C=42pt@R=18pt{
A \ar[rr]^{a\mapsto ea} \ar[dr]_{\text{projection }a\mapsto ea\qquad} && A\\
& eA \ar[ur]_{\text{inclusion}} &
}
\]
So multiplying $e$ means we work inside the direct summand $eA$.
The \defn{corner subalgebra} $eAe=\{eae\mid a\in A\}$ has identity $e$.\footnote{note that this is not a unital subalgebra since the identity element changes; a more proper way to think about this is as a subcategory $\add_{\mod A}(eA)\simeq \add_{\mod eAe}(eAe)$.}
For $A=\prod_{i=1}^r A_i$, the \defn{coordinate idempotents} $e_i=(0,\ldots,1_{A_i},\ldots,0)$ are a complete orthogonal family and are the identities of the individual factors.
In $A=\Bbbk A_n$, the diagonal matrix units $e_{11},\ldots,e_{nn}$ are primitive, since $e_{ii}Ae_{ii}=\Bbbk e_{ii}$.
For $n\geq2$, the sum $e=e_{11}+\cdots+e_{n-1,n-1}$ gives the corner
\[
eAe=\left\{\begin{pmatrix}B&0\\0&0\end{pmatrix}\;\middle|\;B\in\Bbbk A_{n-1}\right\}
\cong\Bbbk A_{n-1}.
\]
\end{remark}

Idempotents also describe maps out of the summands of $A_A$.
\begin{lemma}[Yoneda lemma]\label{lem:idempotent-hom}
For an idempotent $e\in A$ and a right $A$-module $M$, evaluation at $e$ gives a natural (in $M$) vector space isomorphism
\[
\Hom_A(eA,M)\xrightarrow{\sim}Me,\qquad f\longmapsto f(e),
\quad\text{with inverse}\quad m\longmapsto(ea\mapsto ma).
\]
\end{lemma}
\begin{proof}
We have $f(e)=f(e)e$ and $f(ea)=f(e)a$.
If $m=me$, the map $ea\mapsto ma$ is well-defined because $ea=ea'$ implies $ma=mea=mea'=ma'$.
These formulas prove that the maps are inverse; postcomposition by a module map $M\to N$ proves naturality.
\end{proof}




\subsection{Graded algebras}

\begin{definition}\label{def:group-grading}
Let $G$ be a group with identity $e$.
A \defn{$G$-graded $\Bbbk$-algebra} is a $\Bbbk$-algebra $A$ with $1_A\in A_e$ and a decomposition
\[
A=\bigoplus_{g\in G}A_g
\quad\text{such that}\quad
A_gA_h\subset A_{gh}
\quad\text{for all }g,h\in G.
\]
An element of $A_g$ is \defn{homogeneous of degree $g$}.
\end{definition}

\begin{convention}
Unless stated otherwise, a graded algebra is $\Z$-graded and \uline{concentrated in non-negative degrees}.
\end{convention}

Note that the decomposition is into a direct sum, so every element of a graded algebra $A$ can always be written as a \emph{finite} sum of homogeneous elements.

\begin{example}\label{ex:triangular-grading}
For the upper triangular matrix ring $\Bbbk A_n$ of Example~\ref{ex:triangular-algebra}, the decomposition
\[
(\Bbbk A_n)_d=\operatorname{span}_\Bbbk\{e_{ij}\mid j-i=d\}
\]
makes $\Bbbk A_n$ a graded algebra.
\end{example}

\begin{supplementary}
\begin{example}
The algebra $\Bbbk A_3$ also has a $\Z^{\oplus 2}$-grading with
\[
\deg e_{ii}=(0,0),\qquad
\deg e_{12}=(1,0),\qquad
\deg e_{23}=(0,1),\qquad
\deg e_{13}=(1,1).
\]
The homomorphism $\Z^{\oplus 2}\to\Z$, $(p,q)\mapsto p+q$, gives the grading of Example~\ref{ex:triangular-grading}.
Alternatively, setting $\deg e_{ij}=2(j-i)$ gives a $\Z$-grading with $A_1=0$ but $A_2\neq 0$.
\end{example}
\end{supplementary}


\begin{definition}
A \defn{homomorphism of graded algebras} $f:A\to B$ is a $\Bbbk$-algebra homomorphism with $f(A_i)\subset B_i$ for all $i\in\Z$.
\end{definition}

\begin{definition}
A graded algebra $A$ is \defn{locally finite-dimensional} (\defn{lfd}) if $\dim_\Bbbk A_i<\infty$ for every $i\in\Z$.
\end{definition}
%
%\begin{assumption}\label{ass:degree-zero}
%For quadratic and Koszul algebras, $S:=A_0$ is finite-dimensional semisimple.
%\end{assumption}

\subsection{Graded modules}

\begin{definition}\label{def:graded-module}
Let $A=\bigoplus_{g\in G}A_g$ be a $G$-graded algebra.
A \defn{$G$-graded right $A$-module} is a right $A$-module $M$ with a direct sum decomposition
\[
M=\bigoplus_{g\in G}M_g
\quad\text{such that}\quad
M_gA_h\subset M_{gh}
\quad\text{for all }g,h\in G.
\]
An element of $M_g$ is called \defn{homogeneous of degree $g$}.

A $G$-graded $A$-module $M=\bigoplus_{g\in G} M_g$ is \defn{locally finite-dimensional} (\defn{lfd}) if $\dim_{\Bbbk} M_g <\infty$ for all $g\in G$.
\end{definition}

Unless specified otherwise, graded modules are $\Z$-graded; their degrees need not be non-negative.

\phantomsection\label{not:opposite-algebra}
For an algebra $A$, the \defn{opposite algebra} $A^\op$ has the same vector space as $A$, with multiplication $a^\op b^\op=(ba)^\op$.
A left $A$-module is a right $A^\op$-module by $m\cdot a^\op=am$.
Thus $\Mod(A^\op)$ is the category of left $A$-modules and left $A$-linear maps; not to be confused with the opposite category $(\Mod A)^\op$ though.
For a $\Z$-graded algebra $A=\bigoplus_{i\ge 0}A_i$, 
\[
\text{the opposite algebra is also graded with $(A^\op)_i=A_i$.}
\]



\begin{definition}\label{def:shift}
Let $M$ be a graded $A$-module and let $d\in\Z$.
The \defn{degree shift}, or \defn{grading shift}, $M(d)$ is the graded $A$-module whose underlying $A$-module is $M$, but whose grading is
\[
M(d)_j=M_{j+d}.
\]
\end{definition}

By Definition~\ref{def:shift}, $1\in A_0=A(-d)_d$, so $A(-d)$ has a generator in degree $d$.
If $M$ is concentrated in degree $0$, then $M(d)$ is concentrated in degree $-d$.

\begin{example}\label{ex:shifted-row}
For $A=\Bbbk A_3$ with the grading of Example~\ref{ex:triangular-grading}, the module $P_1:=e_{11}A$ consists of matrices supported on the first row.
Its basis $e_{11},e_{12},e_{13}$ has degrees $0,1,2$, respectively; in $P_1(-2)$, the same basis has degrees $2,3,4$.
\end{example}

\begin{definition}
Let $M,N$ be graded $A$-modules.
A homomorphism $f:M\to N$ is \defn{homogeneous of degree $d$} if
$f(M_j)\subset N_{j+d}$ for all $j\in\Z$. A \defn{homomorphism of graded modules} means a homogeneous $A$-module homomorphism of degree $0$, also said to \defn{preserve degrees}.
\end{definition}

\begin{example}\label{ex:matrix-hom}
For $A=\Bbbk A_3$, left multiplication by $e_{12}$ yields a map
\[
\begin{aligned}
e_{22}A&\longrightarrow e_{11}A,\\[-2pt]
\begin{pmatrix}0&0&0\\0&u&v\\0&0&0\end{pmatrix}
&\longmapsto
\begin{pmatrix}0&u&v\\0&0&0\\0&0&0\end{pmatrix}
\qquad(u,v\in\Bbbk).
\end{aligned}
\]
The $u$ and $v$ entries have degrees $0,1$ in the source and $1,2$ in the target, and so the map has degree $1$.
\end{example}

\begin{notation}\label{not:graded-hom}
For graded right $A$-modules $M,N$, the space of homomorphisms preserving degrees is
\begin{equation}\label{eq:graded-hom}
\Hom_A^{\Z}(M,N)=\{f\in\Hom_A(M,N)\mid f(M_j)\subset N_j\text{ for all }j\in\Z\}.
\end{equation}
Thus $\Hom_A^{\Z}(M,N(d))$ consists of homomorphisms $M\to N$ of degree $d$, since $N(d)_j=N_{j+d}$.
Note that, like the grading shift functor, different literature has different notation for this (there does not seem to be any widely agreed notation).
\end{notation}

\phantomsection\label{not:graded-module-categories}
We write 
\begin{align*}
\Grmod A &:= \text{category of graded right $A$-modules with morphism spaces $\Hom_A^{\Z}(M,N)$},\\
\grmod A &:= \text{full subcategory of $\Grmod A$ consisting of locally finite-dimensional modules}.
\end{align*}
Thus $\Grmod(A^\op)$ is the category of graded left $A$-modules and degree-preserving maps.
Forgetting the grading gives a functor $\Grmod A\to\Mod A$.

\begin{remark}
$\grmod A$ is somewhat analogous to $\mod A$ but not really the same:
If $A$ is locally finite-dimensional, every finitely generated graded $A$-module $M$ is locally finite-dimensional. ($\because$ a finite homogeneous generating set of a graded $A$-module $M$ gives a surjection from a finite direct sum of shifts of $A$ to $M$, so every $M_i$ has to be finite-dimensional.)  However, a locally finite-dimensional $M$ need not be finitely generated.
\end{remark}
 
For $d\in\Z$, the shift of Definition~\ref{def:shift} is a functor \[(d):\Grmod A\to\Grmod A.\]
It sends $f:M\to N$ to the same underlying map $f(d):M(d)\to N(d)$.
Since $M(d)(-d)=M$ and $f(d)(-d)=f$, the shift is an equivalence with inverse $(-d)$; it restricts to an equivalence $\grmod A\to\grmod A$.

\begin{remark}\label{rem:graded-exactness}
The notion of exactness makes sense for graded module categories -- a sequence in $\Grmod A$ is exact if the underlying sequence in $\Mod A$ is exact.
In particular, we talk about functors invovle such categories being exact.
\begin{supplementary}
For $L,M,N\in\Grmod A$ with maps $i:L\to M$ and $p:M\to N$, the sequence $0\to L\xrightarrow{i}M\xrightarrow{p}N\to0$ is exact if and only if
\[
0\longrightarrow L_j\xrightarrow{i|_{L_j}}M_j\xrightarrow{p|_{M_j}}N_j\longrightarrow0
\quad\text{is exact for every }j\in\Z.
\]
Indeed, $i$ and $p$ preserve the direct sum decompositions, so their kernels and images decompose into the kernels and images of the restrictions.
\end{supplementary}
\end{remark}

\begin{example}\label{ex:short-exact-row}
Let $A=\Bbbk A_2$. 
Consider $P_1:=e_{11}A$ which has $\rad P_1 = \Bbbk e_{12}$ and $L_1:=P_1/\rad P_1$.
Then we have a short exact sequence in $\Grmod A$:
\[
0\longrightarrow \rad P_1\xrightarrow{i} P_1
\xrightarrow{p}P_1/\rad P_1\longrightarrow0,
\qquad
i(c e_{12})=c e_{12},\quad p(m)=\overline m.
\]
Note that $\deg \rad P_1 = 1$ and $\deg L_1=0$.  Let $L_2:=\Bbbk e_{22} = e_{22} A$, which we treat as a graded $A$-module concentrated in degree $0$.  Note that $L_2 \cong \rad P_1$ as ungraded modules but we only have $L_2\cong (\rad P_1) (1)$, which yields the following short exact sequence in $\Grmod A$:
\[
0\longrightarrow L_2(-1)\xrightarrow{} P_1
\xrightarrow{}P_1/\rad P_1\longrightarrow0,
\]
\end{example}

\begin{lemma}\label{lem:homogeneous-hom-sum}
For graded right $A$-modules $M,N$, there is an inclusion 
\[
\bigoplus_{d\in\Z}\Hom_A^{\Z}(M,N(d))\subset\Hom_A(M,N), \text{ which is equality when $M$ is finitely generated.}
\]
\end{lemma}

\begin{proof}
For $f\in\Hom_A(M,N)$, define $f_d(m)=f(m)_{j+d}$ when $m\in M_j$, and extend linearly.
For $a\in A_i$, we have $f_d(ma)=f(ma)_{j+i+d}=f(m)_{j+d}a$, so $f_d$ is $A$-linear of degree $d$.
The sum of the homogeneous Hom spaces is direct: on $M_j$, distinct degrees have values in distinct summands of $N$.

If $M$ is finitely generated, the homogeneous components of a finite generating set form a finite homogeneous generating set.
Only finitely many $f_d$ are non-zero on these homogeneous generators, hence on $M$, giving $f=\sum_d f_d$.
\end{proof}

\subsection{Homogeneous ideals and quotients}

\begin{definition}
Let $A=\bigoplus_{i\in\Z}A_i$ be a graded algebra.
A right ideal $I\subset A$ is \defn{homogeneous} if
\[
I=\bigoplus_{i\in\Z}(I\cap A_i).
\]
Equivalently, whenever $a\in I$ and $a=a_0+\cdots+a_N$ is its decomposition into homogeneous components, each $a_i$ also belongs to $I$.
\end{definition}

For quotient algebras we need two-sided ideals.
If $I$ is a homogeneous two-sided ideal, then
\[
A/I=\bigoplus_{i\in\Z} A_i/(I\cap A_i)
\]
is again a graded algebra.

\begin{example}\label{ex:quadratic-matrix-quotient}
In $A=\Bbbk A_3$ with the grading of Example~\ref{ex:triangular-grading}, the two-sided ideal $I=(e_{13})=\Bbbk e_{13}$ is homogeneous.
Writing $\overline e_{ij}$ for the image of $e_{ij}$,
\[
(A/I)_0\cong\Bbbk^{\oplus 3},\qquad
(A/I)_1=\Bbbk\overline e_{12}\oplus\Bbbk\overline e_{23},\qquad
(A/I)_d=0\quad(d\geq 2).
\]
In particular, $\overline e_{12}\overline e_{23}=0$.
\end{example}

\begin{example}\label{ex:nonhomogeneous-ideal}
Let $A=\Bbbk[x]$ with $\deg x=1$, and let $I=(x-x^2)$.
The ideal $I$ is two-sided because $A$ is commutative.
Every element of $I$ vanishes at $x=1$, whereas neither $x$ nor $x^2$ does.
Thus $x-x^2\in I$ but $x,x^2\notin I$, so $I$ is not homogeneous.
Writing $\overline f=f+I$, we have $0\neq\overline x=\overline{x^2}$.
The images of $A_1$ and $A_2$ therefore intersect non-trivially, so $A/I$ does not inherit the quotient grading.
As an ungraded algebra, the quotient is semisimple:
\[
\Bbbk[x]/(x-x^2)\xrightarrow{\sim}\Bbbk\oplus\Bbbk,
\qquad \overline f\longmapsto(f(0),f(1)).
\]
The inverse sends $(a,b)$ to $a+(b-a)\overline x$.
\end{example}

\begin{definition}\label{def:positive-ideal}
The \defn{positive degree ideal} of a graded algebra $A$ is $A_+:=\bigoplus_{i>0}A_i$.\footnote{It is also called the augmentation ideal; we avoid that terminology.}
\end{definition}

The subspace $A_+$ is a homogeneous two-sided ideal.

\subsection{Tensor products of vector spaces}

\begin{definition}\label{def:vector-tensor}
Choose a basis $\{u_i\mid i\in I\}$ of $U$ and a basis $\{v_j\mid j\in J\}$ of $V$.
The \defn{tensor product} $U\otimes_\Bbbk V$ is the vector space with basis $\{u_i\otimes v_j\mid i\in I,\ j\in J\}$.
\end{definition}

The assignment $(u_i,v_j)\mapsto u_i\otimes v_j$ extends uniquely to a bilinear map $U\times V\to U\otimes_\Bbbk V$ sending $(u,v)\mapsto u\otimes v$.  Note that this means that elements like $(au_1+bu_2)\otimes (cv_1+dv_2)$ in $U\otimes_\Bbbk V$ is by definition the linear combination $ac(u_1\otimes v_1) + ad(u_1\otimes v_2) + bc(u_2\otimes v_1) + bd(u_2\otimes v_2)$.


% An element of the form $u\otimes v$ is called a \defn{pure tensor}.
% \comm{do not introduce useless terminology!}
\begin{proposition}[Universal property]\label{prop:vector-tensor-universal}
Every bilinear map $b:U\times V\to W$ has a unique linear factorisation (commutative diagram)
\[
\xymatrix@C=60pt@R=22pt{
U\times V \ar[r]^{(u,v)\mapsto u\otimes v} \ar[dr]_{\forall \text{ bilinear } b\phantom{...}}
& U\otimes_\Bbbk V \ar@{-->}[d]^{\exists!\,\widetilde b}\\
& W.
}
\]
% File-specific editorial rule for 2026_Koszul.tex: preserve the author's
% manual spacing adjustments, including \phantom{...} in diagram labels.
% Do not remove or change those adjustments unless explicitly requested.
Consequently, $U\otimes_\Bbbk V$ is independent of the chosen bases up to a unique isomorphism preserving the bilinear map.
\end{proposition}


\begin{supplementary}
\begin{proof}
The formula $\widetilde b(u_i\otimes v_j)=b(u_i,v_j)$ determines a unique linear map $\widetilde b:U\otimes_\Bbbk V\to W$.
Bilinearity gives $\widetilde b(u\otimes v)=b(u,v)$ for all $u\in U$ and $v\in V$.

Suppose $t:U\times V\to T$ and $t':U\times V\to T'$ both satisfy the universal property.
They give unique linear maps $a:T\to T'$ and $b:T'\to T$ with $a\circ t=t'$ and $b\circ t'=t$.
Both $b\circ a:T\to T$ and $\id_T:T\to T$ factor $t$ through itself, since
\[
(b\circ a)\circ t=b\circ t'=t=\id_T\circ t.
\]
Uniqueness for $(T,t)$ gives $b\circ a=\id_T$.
Likewise, $(a\circ b)\circ t'=a\circ t=t'=\id_{T'}\circ t'$, so uniqueness for $(T',t')$ gives $a\circ b=\id_{T'}$.
\end{proof}
\end{supplementary}

For finite bases with $|I|=r$ and $|J|=s$, we obtain $\dim_\Bbbk(U\otimes_\Bbbk V)=rs$.
We use the canonical identifications
\[
\begin{aligned}
(U\otimes_\Bbbk V)\otimes_\Bbbk W&\xrightarrow{\sim}U\otimes_\Bbbk(V\otimes_\Bbbk W),
& (u\otimes v)\otimes w&\longmapsto u\otimes(v\otimes w),\\
\Bbbk\otimes_\Bbbk U&\xrightarrow{\sim}U,
& c\otimes u&\longmapsto cu,
\end{aligned}
\]
and set $U^{\otimes_\Bbbk 0}:=\Bbbk$.


\subsection{Direct sums and direct products}\label{subsec:sums-products}

For an arbitrary family $(M_i)_{i\in I}$ of right $A$-modules, indexed by a set $I$, the \defn{direct product} and \defn{direct sum} are
\[
\begin{aligned}
\prod_{i\in I}M_i&=\{(m_i)_{i\in I}\mid m_i\in M_i\text{ for all }i\in I\},\\
\bigoplus_{i\in I}M_i&=\{(m_i)_{i\in I}\in\prod_{i\in I}M_i\mid m_i=0\text{ for all but finitely many }i\}.
\end{aligned}
\]
Addition and the $A$-action are componentwise; the same definitions apply to vector spaces by taking $A=\Bbbk$.
Thus $\bigoplus_{i\in I} M_i\subset\prod_{i\in I} M_i$, with equality for finite $I$.
For example, $(1,1,\ldots)\in\prod_{i\geq0}\Bbbk$ does not belong to $\bigoplus_{i\geq0}\Bbbk$.
In a direct sum, an element is a finite sum of its components; no notion of convergence is involved.
For instance, ordinary polynomials use non-negative powers, whereas Laurent polynomials allow integer powers, but both are finite sums:
\[
\Bbbk[x]=\bigoplus_{n\geq0}\Bbbk x^n,\qquad
\Bbbk[x,x^{-1}]=\bigoplus_{n\in\Z}\Bbbk x^n.
\]
\phantomsection\label{not:formal-power-series}
The \defn{formal power series ring} allows arbitrary coefficients:
\[
\Bbbk[[x]]=\left\{\sum_{n\geq0}a_nx^n\;\middle|\;a_n\in\Bbbk\text{ for every }n\geq0\right\}
\cong\prod_{n\geq0}\Bbbk x^n
\quad\text{as vector spaces}.
\]
\begin{supplementary}
Addition is coefficientwise, and multiplication is
\[
\Bigl(\sum_{n\geq0}a_nx^n\Bigr)\Bigl(\sum_{n\geq0}b_nx^n\Bigr)
=\sum_{n\geq0}\Bigl(\sum_{j=0}^n a_jb_{n-j}\Bigr)x^n.
\]
\end{supplementary}


Tensor products distribute over direct sums in either variable.
For vector spaces $U,V$ and arbitrary families $(U_i)_{i\in I}$, $(V_j)_{j\in J}$ indexed by sets $I,J$, there are natural linear isomorphisms
\begin{equation}\label{eq:vector-tensor-sums}
\begin{aligned}
\Bigl(\bigoplus_{i\in I}U_i\Bigr)\otimes_\Bbbk V
&\xrightarrow{\sim}\bigoplus_{i\in I}(U_i\otimes_\Bbbk V),
& (u_i)_{i\in I}\otimes v&\longmapsto(u_i\otimes v)_{i\in I},\\
U\otimes_\Bbbk\Bigl(\bigoplus_{j\in J}V_j\Bigr)
&\xrightarrow{\sim}\bigoplus_{j\in J}(U\otimes_\Bbbk V_j),
& u\otimes(v_j)_{j\in J}&\longmapsto(u\otimes v_j)_{j\in J}.
\end{aligned}
\end{equation}
The displayed bilinear assignments induce these linear maps by Proposition~\ref{prop:vector-tensor-universal}.
On the $i$th summand, the inverse of the first map sends $u_i\otimes v$ to the element with $u_i$ in coordinate $i$, tensored with $v$.
On the $j$th summand, the inverse of the second sends $u\otimes v_j$ to $u$ tensored with the element having $v_j$ in coordinate $j$.
Extending over the direct sums gives the two inverses.

\begin{advanced}
For direct products, the corresponding map need not be an isomorphism.

\begin{lemma}\label{lem:vector-tensor-products}
For an arbitrary family $(U_i)_{i\in I}$ of vector spaces, indexed by a set $I$, and a vector space $V$, the coordinate projections define a linear map
\begin{equation}\label{eq:vector-tensor-products}
\Bigl(\prod_{i\in I}U_i\Bigr)\otimes_\Bbbk V
\longrightarrow\prod_{i\in I}(U_i\otimes_\Bbbk V),
\qquad (u_i)_i\otimes v\longmapsto(u_i\otimes v)_i.
\end{equation}
This map is injective, and is an isomorphism if $I$ is finite or $V$ is finite-dimensional.
The same statements hold with the two tensor factors reversed.
\end{lemma}

\begin{proof}
The bilinear map $\bigl(\prod_{i\in I}U_i\bigr)\times V\to\prod_{i\in I}(U_i\otimes_\Bbbk V)$, $((u_i)_i,v)\mapsto(u_i\otimes v)_i$, induces \eqref{eq:vector-tensor-products} by Proposition~\ref{prop:vector-tensor-universal}.
Every element can be written $\sum_{\ell=1}^r u^{(\ell)}\otimes v_\ell$ with the $v_\ell$ linearly independent: choose a basis of the span of its finitely many second factors.
If its image is zero, then $\sum_{\ell=1}^r u_i^{(\ell)}\otimes v_\ell=0$ for every $i$.
The tensor-product basis construction gives $u_i^{(\ell)}=0$ for all $i,\ell$, proving injectivity.
If $v_1,\ldots,v_r$ is a basis of $V$, the inverse is
\[
\prod_{i\in I}(U_i\otimes_\Bbbk V)\longrightarrow
\Bigl(\prod_{i\in I}U_i\Bigr)\otimes_\Bbbk V,\qquad
\Bigl(\sum_{\ell=1}^r u_i^{(\ell)}\otimes v_\ell\Bigr)_i
\longmapsto\sum_{\ell=1}^r(u_i^{(\ell)})_i\otimes v_\ell.
\]
For finite $I$, products equal sums.
Reversing the tensor factors gives the other statements, since $U\otimes_\Bbbk V\to V\otimes_\Bbbk U$, $u\otimes v\mapsto v\otimes u$, is an isomorphism.
\end{proof}

\begin{example}\label{ex:tensor-product-product-failure}
The comparison in \eqref{eq:vector-tensor-products} need not be surjective when $V$ is infinite-dimensional.
Take $U_i=\Bbbk$ for $i\geq1$ and $V=\bigoplus_{j\geq1}\Bbbk v_j$.
Under $\Bbbk\otimes_\Bbbk V\cong V$, the image of the product comparison in \eqref{eq:vector-tensor-products} is
\[
\left\{(w_i)_i\in\prod_{i\geq1}V\;\middle|\;
\dim_\Bbbk\operatorname{span}_\Bbbk\{w_i\mid i\geq1\}<\infty\right\}.
\]
Indeed, the image of $\sum_{\ell=1}^r(c_i^{(\ell)})_i\otimes w_\ell$ has every entry in $\operatorname{span}_\Bbbk\{w_1,\ldots,w_r\}$.
Conversely, expand all entries in a basis of their common finite-dimensional span.
Since the $v_i$ are linearly independent, $(v_i)_i$ is not in the image.
\end{example}
\end{advanced}

\subsection{Bimodules and enveloping algebras}


\begin{definition}\label{def:bimodule}
For $\Bbbk$-algebras $B,C$, a \defn{bimodule} $ {}_BM_C$ is a vector space with $\Bbbk$-bilinear left $B$- and right $C$-module actions satisfying $(bm)c=b(mc)$ for all $b\in B$, $c\in C$, and $m\in M$.
The subscripts indicate the left and right acting algebras; when $B=C=A$, we call $M$ an \defn{$A$-bimodule}.
A \defn{bimodule homomorphism} $f:{}_BM_C\to{}_BN_C$ is a linear map satisfying $f(bmc)=bf(m)c$ for all $b\in B$, $c\in C$, and $m\in M$.
For graded $B,C$, a bimodule $M$ is \defn{graded} if $M=\bigoplus_{i\in\Z}M_i$ and $B_iM_jC_k\subset M_{i+j+k}$ for all $i,j,k\in\Z$.
A \defn{graded bimodule homomorphism} is a bimodule homomorphism $f:M\to N$ satisfying $f(M_i)\subset N_i$ for all $i\in\Z$.
\end{definition}

\phantomsection\label{not:bimodule-notation}
\begin{convention}
\begin{itemize}
\item Bilinearity implies that $\Bbbk$ acts centrally: $(c1_B)m=cm=m(c1_C)$ for $c\in\Bbbk$.  
From the perspective of general module theory over general rings (algebras over non-field commutative rings), we should have said that our bimodules are $\Bbbk$-central bimodule.  Since we alawys work over a field, all bimodules we considered are always $\Bbbk$-central and we will not write such adjective every time.

\item  Instead of saying $(B,C)$-bimodule, we usually just write ${}_BM_C$.
It would be natural to write $B\Mod C$ for the category of $(B,C)$-bimodules, but since $\Bbbk$-centrality comes for free from the assumption of $\Bbbk$ being a field, we can use $\Mod C\otimes_\Bbbk B^\op$ instead.
\end{itemize}
\end{convention}

\begin{example}
\begin{itemize}
\item Multiplication makes $A_A$ a graded right module and ${}_AA_A$ a graded bimodule; each degree component is a bimodule ${}_{A_0}(A_i)_{A_0}$.

\item For a bimodule ${}_AM_A$ and idempotents $e,f\in A$, the space $eMf:=\{enf\mid n\in M\}$ is a bimodule ${}_{eAe}(eMf)_{fAf}$.
\end{itemize}
\end{example}


\begin{definition}\label{def:enveloping-algebra}
For a $\Bbbk$-algebra $A$, its \defn{enveloping algebra} is $A^\en:=A\otimes_\Bbbk A^\op$, with multiplication extended bilinearly from
\[
(a\otimes b^\op)(a'\otimes(b')^\op)=aa'\otimes(b'b)^\op,
\qquad 1_{A^\en}=1_A\otimes1_A^\op
\quad(a,b,a',b'\in A).
\]
\end{definition}

A bimodule ${}_AM_A$ becomes a right $A^\en$-module by
\begin{equation}\label{eq:enveloping-action}
m\cdot(a\otimes b^\op)=bma
\quad(a,b\in A,\ m\in M).
\end{equation}
\begin{supplementary}
The first tensor factor records the right action, and the second the left action:
\[
\bigl(m\cdot(a\otimes b^\op)\bigr)\cdot(a'\otimes(b')^\op)
=b'bmaa'
=m\cdot\bigl(aa'\otimes(b'b)^\op\bigr).
\]
Conversely, for $M\in\Mod A^\en$, define $ma=m\cdot(a\otimes1)$ and $bm=m\cdot(1\otimes b^\op)$.
The two actions commute because $a\otimes1$ and $1\otimes b^\op$ commute in $A^\en$.
For bimodules ${}_AM_A$, ${}_AN_A$, a linear map $u:M\to N$ satisfies
\[
u(bma)=bu(m)a\quad(a,b\in A,\ m\in M)
\quad\Longleftrightarrow\quad
u\in\Hom_{A^\en}(M,N).
\]
Thus $A\Mod A$ is identified with $\Mod A^\en$.
More generally, for $\Bbbk$-algebras $B,C$, the vector space $C\otimes_\Bbbk B^\op$ is an algebra with
\[
(c\otimes b^\op)(c'\otimes(b')^\op)=cc'\otimes(b'b)^\op,
\qquad 1=1_C\otimes1_B^\op.
\]
The same construction identifies the categories
\[
B\Mod C\simeq\Mod(C\otimes_\Bbbk B^\op),
\qquad m\cdot(c\otimes b^\op)=bmc.
\]
\end{supplementary}

\begin{advanced}
Grant--Iyama \cite[\S2, p.~2591]{GI20} identify bimodules with left modules over $A^\en$, using $(a\otimes b^\op)m=amb$.
Our right module convention is \eqref{eq:enveloping-action}: $m\cdot(a\otimes b^\op)=bma$.
Both describe the same bimodules; the side on which $A^\en$ acts determines the formula.

\begin{remark}\label{rem:enveloping-semisimple}
For a finite-dimensional semisimple $\Bbbk$-algebra $S$, semisimplicity of $S^\en$ requires an additional condition.
For a field extension $K/\Bbbk$, give $K\otimes_\Bbbk S$ the multiplication $(c\otimes s)(d\otimes t)=cd\otimes st$.
The algebra $S$ is called \defn{separable over $\Bbbk$} if $K\otimes_\Bbbk S$ is semisimple for every field extension $K/\Bbbk$.
By \cite[Theorem~6.1.2, p.~105]{DK94},
\[
S^\en\text{ is semisimple}
\quad\Longleftrightarrow\quad
S\text{ is separable over }\Bbbk.
\]
For a product of matrix algebras, every field extension $K/\Bbbk$ gives
\[
K\otimes_\Bbbk\prod_{j=1}^n\Mat_{n_j}(\Bbbk)
\xrightarrow{\sim}\prod_{j=1}^n\Mat_{n_j}(K),
\qquad c\otimes(a_j)_j\longmapsto(ca_j)_j.
\]
The target is semisimple by Theorem~\ref{thm:artin-wedderburn}.
Thus products of matrix algebras over $\Bbbk$ are separable; in particular, every finite-dimensional semisimple algebra over an algebraically closed field is separable.
Semisimplicity of $S$ alone does not imply semisimplicity of $S^\en$ over an arbitrary field.
\end{remark}

\begin{example}\label{ex:inseparable-envelope}
Let $\Bbbk=\mathbb F_2(t)$, where $t$ is an indeterminate, and let $S=\Bbbk[u]/(u^2-t)$, also writing $u$ for its residue class.
The element $t$ is not a square in $\Bbbk$: an equality $p(t)^2=tq(t)^2$ for non-zero polynomials would give $2\deg p=1+2\deg q$.
Thus $u^2-t$ is irreducible and $S$ is a field, hence semisimple.
However, since $S$ is commutative and has characteristic $2$,
\[
S^\en=S\otimes_\Bbbk S
\cong S[v]/(v^2-t)
=S[v]/((v-u)^2)
\cong S[\varepsilon]/(\varepsilon^2).
\]
Here the first tensor factor supplies the coefficients, $v$ corresponds to $1\otimes u$, and $\varepsilon=1\otimes u-u\otimes1$.
The ideal $S\varepsilon$ is non-zero, has square zero, and is the unique maximal ideal, so $\rad(S^\en)=S\varepsilon\neq0$.
Thus $S^\en$ is not semisimple, and the field extension $S/\Bbbk$ already shows that $S$ is not separable over $\Bbbk$.
\end{example}
\end{advanced}


%%%%%%%%%%%%%%%%%%%%%%%%%%%%%%%%%%%%%%%%%%%%%%%%%%%%%%%%%%%%%%%%%%

\subsection{\texorpdfstring{Graded $\Bbbk$-linear duals}{Graded k-linear duals}}

A linear functional on a graded vector space can be non-zero on infinitely many degrees.
To obtain a graded space, retain only finite sums of homogeneous functionals.

\begin{definition}\label{def:linear-duals}
For a graded vector space $V=\bigoplus_{i\in\Z}V_i$, the \defn{graded $\Bbbk$-linear dual} is the graded vector space
\[
\bbD(V)=\bigoplus_{i\in\Z}\bbD(V)_i,
\qquad \bbD(V)_i=D(V_{-i})=\Hom_\Bbbk^\Z(V,\Bbbk(i)).
\]
Note that here $\Bbbk$ is regarded as a graded module concentrated in degree $0$.
\end{definition}

Linear duality reverses the side of a module action.
For a right $A$-module $M$ and a left $A$-module $N$, set
\begin{equation}\label{eq:linear-dual-actions}
(af)(m)=f(ma),\qquad (ga)(n)=g(an)
\quad\text{for }a\in A,\ f\in D(M),\ g\in D(N),\ m\in M,\ n\in N.
\end{equation}
For graded modules these actions restrict to the graded linear duals: if $a\in A_p$ and $f\in\bbD(M)_q$, then $af$ vanishes outside $M_{-p-q}$.


\begin{example}\label{ex:row-linear-dual}
Let $A=\Bbbk A_3$ and $P_1=e_{11}A$ as in Example~\ref{ex:shifted-row}.
The functionals $\varepsilon_{1j}(e_{1\ell})=\delta_{j\ell}$ for $1\leq j,\ell\leq 3$ form a basis of $\bbD(P_1)$, with
\[
\deg\varepsilon_{1j}=1-j,\qquad
e_{12}\varepsilon_{12}=\varepsilon_{11},\qquad
e_{23}\varepsilon_{13}=\varepsilon_{12},\qquad
e_{jj}\varepsilon_{1j}=\varepsilon_{1j}.
\]
The actions follow from \eqref{eq:linear-dual-actions}.
%For comparison with an infinite-dimensional space, take $V=\bigoplus_{i\geq 0}\Bbbk v_i$ with $\deg v_i=i$.
%The functional $\lambda(v_i)=1$ for every $i$ belongs to $D(V)$ but not to $\bbD(V)$.
\end{example}


For any $u\in \Hom_A(M,N)$, we have a dual map
\[D(u):D(N)\to D(M), \qquad f\mapsto f\circ u,
\] 
which is a left $A$-linear map, i.e. $D(u)\in \Hom_{A^\op}(D(N), D(M))$.
This gives rise to four exact (contravariant) functors (without any finiteness assumption):
\begin{equation}\label{eq:linear-dual-functors}
\begin{aligned}
D:(\Mod A)^\op&\longrightarrow\Mod(A^\op),
&D:(\Mod(A^\op))^\op&\longrightarrow\Mod A,\\
\bbD:(\Grmod A)^\op&\longrightarrow\Grmod(A^\op),
&\bbD:(\Grmod(A^\op))^\op&\longrightarrow\Grmod A.
\end{aligned}
\end{equation}
Exactness follows from Example~\ref{ex:linear-duality}, degree by degree for $\bbD$.

For these to be duality, we need to recover a module from its bidual; this requires some finiteness conditions.

\begin{lemma}\label{lem:linear-bidual}
Let $A$ be a $\Bbbk$-algebra.
\begin{enumerate}
\item The functor $D$ defines a duality between finite-dimensional right and left $A$-modules.  In particular, if $A$ is finite-dimensional, then we have a duality on finitely generated modules
\[
D: (\mod A)^\op \xrightarrow{\sim}\mod (A^\op), \qquad D: (\mod (A^\op))^\op \xrightarrow{\sim}\mod A.
\]

\item Suppose that $A$ is graded.  For locally finite-dimensional $M\in \grmod A$, evaluation is a natural graded module isomorphism
\begin{equation}\label{eq:linear-bidual-evaluation}
\mathrm{ev}_M:M\xrightarrow{\sim}\bbD(\bbD(M)),\qquad
m\longmapsto(f\mapsto f(m)).
\end{equation}
In particular, we have a duality between locally finite-dimensional graded right and left $A$-modules:
\begin{equation}\label{eq:graded-linear-dual}
\bbD:(\grmod A)^\op\xrightarrow{\sim}\grmod(A^\op),\qquad
\bbD:(\grmod(A^\op))^\op\xrightarrow{\sim}\grmod A.
\end{equation}
\end{enumerate}
\end{lemma}

\begin{proof}
First we show that $\mathrm{ev}_M$ is $A$-linear.
Indeed, 
\[
\mathrm{ev}_M(ma)(f)=f(ma)
=\mathrm{ev}_M(m)(af)
=(\mathrm{ev}_M(m)a)(f)
\quad(m\in M,\ a\in A,\ f\in D(M)).
\]

For (1), the claim follows from Example~\ref{ex:linear-duality}.
For (2), apply (1) degree-wise.
\end{proof}

\begin{advanced}
For a locally finite-dimensional graded vector space $V$, restriction gives
\[
\bbD(V)=\bigoplus_{i\in\Z}D(V_{-i})\subset
D(V)\cong\prod_{i\in\Z}D(V_i).
\]
The inclusion is equality exactly when only finitely many $V_i$ are non-zero.
For an arbitrary vector space, evaluation $V\to D(D(V))$ is injective but need not be surjective; hence the functors in \eqref{eq:linear-dual-functors} are not dualities on all modules.

\begin{example}
Graded linear duality need not preserve finite generation, even for a polynomial ring.
Let $A=\Bbbk[x]$ with $\deg x=1$ and let $\varepsilon_n(x^m)=\delta_{nm}$ for $n,m\geq0$.
Then $\deg\varepsilon_n=-n$, $x\varepsilon_0=0$, and $x\varepsilon_n=\varepsilon_{n-1}$ for $n\geq1$.
Any finite set of elements of $\bbD(A)$ lies in the span of $\varepsilon_0,\ldots,\varepsilon_N$ for some $N$, and multiplication by $x$ preserves that span.
Thus $\bbD(A)$ is locally finite-dimensional but not finitely generated.
\end{example}
\end{advanced}

For vector spaces $V,W$, a functional $f\in D(V)$ and a vector $w\in W$ define a linear map $V\to W$, $v\mapsto f(v)w$.
Finite sums of these maps give the following description of Hom.

\begin{proposition}\label{prop:linear-Hom-tensor}
For vector spaces $V,W$, there is a natural linear map
\begin{equation}\label{eq:linear-Hom-tensor}
W\otimes_\Bbbk D(V)\longrightarrow\Hom_\Bbbk(V,W),
\qquad w\otimes f\longmapsto\bigl(v\mapsto f(v)w\bigr).
\end{equation}
It is an isomorphism if either $V$ or $W$ is finite-dimensional.
\end{proposition}

\begin{supplementary}
\begin{proof}
If $V$ has a finite basis $\{v_i\mid 1\leq i\leq n\}$ with dual basis $\{v_i^*\mid 1\leq i\leq n\}$, the inverse of \eqref{eq:linear-Hom-tensor} is
\[
\Hom_\Bbbk(V,W)\longrightarrow W\otimes_\Bbbk D(V),
\qquad F\longmapsto\sum_{i=1}^nF(v_i)\otimes v_i^*,
\]
since $v=\sum_{i=1}^n v_i^*(v)v_i$.
If instead $W$ has a finite basis $\{w_j\mid 1\leq j\leq m\}$, write $F(v)=\sum_{j=1}^m f_j(v)w_j$; the uniquely determined coordinate functionals $f_j\in D(V)$ give the inverse $F\mapsto\sum_{j=1}^m w_j\otimes f_j$.
The formulas commute with precomposition and postcomposition, so the comparison is natural.
\end{proof}
\end{supplementary}

\begin{assumption}\label{ass:basic-semisimple}
For simplicity, from now on and unless stated otherwise, we assume the degree $0$ part $A_0$ of a graded algebra to be semisimple of the form
\[
S=\bigoplus_{i=1}^n\Bbbk e_i
\]
with coordinate idempotents $e_1,\ldots,e_n$.
\end{assumption}

The coordinate idempotents are the primitive idempotents of $S$, so they involve no choice.
Statements already proved for an arbitrary finite-dimensional semisimple $S$ keep their generality.
A remark headed \emph{General $S$} marks each place where an arbitrary finite-dimensional semisimple $S$ needs extra care; usually this means replacing the $\Bbbk$-linear dual by the one-sided $S$-duals of the advanced material below, or choosing a symmetrising form.
%To dualise these components while retaining both actions, we distinguish maps to $S$ that are right or left $S$-linear.
%For duality of arbitrary $A$-modules we already have the canonical functors in \eqref{eq:linear-dual-functors}.

For a bimodule ${}_SM_S$, the $\Bbbk$-linear dual retains both actions:
\begin{equation}\label{eq:linear-bimodule-dual-action}
(s\lambda t)(m)=\lambda(tms)\qquad(s,t\in S,\ \lambda\in D(M),\ m\in M).
\end{equation}
With precomposition on morphisms, we obtain functors
\[
D:(\Mod S^\en)^\op\longrightarrow\Mod S^\en,\qquad
\bbD:(\Grmod S^\en)^\op\longrightarrow\Grmod S^\en,
\]
whose restrictions to finite-dimensional bimodules, and to locally finite-dimensional graded bimodules, are dualities by Lemma~\ref{lem:linear-bidual}.
\begin{supplementary}
Evaluation preserves both actions: for $m\in M$, $\lambda\in D(M)$, and $s,t\in S$,
\[
(s\,\mathrm{ev}_M(m)\,t)(\lambda)=\mathrm{ev}_M(m)(t\lambda s)=(t\lambda s)(m)=\lambda(smt)=\mathrm{ev}_M(smt)(\lambda).
\]
Applying this in each degree gives the graded statement.
\end{supplementary}

Dualising reverses the two idempotents: for $1\leq i,j\leq n$,
\begin{equation}\label{eq:dual-blocks}
e_jD(M)e_i=\{\lambda\in D(M)\mid\lambda(e_kMe_\ell)=0\text{ unless }(k,\ell)=(i,j)\}\cong D(e_iMe_j),
\end{equation}
since $(e_j\lambda e_i)(m)=\lambda(e_ime_j)$ by \eqref{eq:linear-bimodule-dual-action}.


\begin{advanced}[suppressfirstparskip=true,innertopmargin=0pt]
\subsection{$S$-dual}

Under Assumption~\ref{ass:basic-semisimple}, the $\Bbbk$-linear dual is all we need.
For an arbitrary finite-dimensional semisimple $S$, the choice-free substitutes are the one-sided $S$-duals below; they are the duals used in \cite[\S2.7]{BGS96}.
Under Assumption~\ref{ass:basic-semisimple} they are canonically isomorphic to $D$; see Remark~\ref{rem:standard-form}.

\emph{One-sided $S$-duals.}

\begin{definition}\label{def:S-duals}
For a right $S$-module $M$ and a left $S$-module $N$, define their \defn{$S$-duals} by
\[
\begin{aligned}
M^*&:=\Hom_S(M,S)\in \Mod(S^\op),&
(sf)(m)&:=sf(m),\\
{}^*N&:=\Hom_{S^\op}(N,S) \in \Mod S,&
(gs)(n)&:=g(n)s.
\end{aligned}
\]
Here $s\in S$, $m\in M$, $n\in N$, $f\in M^*$, and $g\in{}^*N$.
\end{definition}

%The position of the star records the side of the original action: $f(ms)=f(m)s$ in $M^*$, whereas $g(sn)=sg(n)$ in $ {}^*N$.
%\comm{why repeat what is already written????}

Similar to $\Bbbk$-lienar duality, this definition extends to exact contravariant functors: for a right module map $u:M\to N$, define $u^*:N^*\to M^*$ by precomposition $u^*(f)=f\circ u$; for a left module map $u:M\to N$, the same recipe defines ${}^*u:{}^*N\to{}^*M$.
This yields the desired functors
\[
(-)^*:(\Mod S)^\op\longrightarrow\Mod(S^\op),
\qquad
{}^*(-):(\Mod(S^\op))^\op\longrightarrow\Mod S.
\]

We also have graded analogue similar to Definition~\ref{def:linear-duals}.

\begin{definition}\label{def:graded-S-duals}
For graded right and left $S$-modules $M$, respectively, define their \defn{graded $S$-duals} by
\[
M^{\grstar}:=\bigoplus_{i\in\Z}(M_{-i})^*
=\bigoplus_{i\in\Z}\Hom_S^{\Z}(M,S(i)),\qquad
{}^{\grstar}M:=\bigoplus_{i\in\Z}{}^*(M_{-i})=\bigoplus_{i\in\Z}\Hom_{S^\op}^{\Z}(M,S(i)),
\]

where each indicated summand has degree $i$.
%The shift is $S(i)$, since $S(i)_{-i}=S_0$.
\end{definition}

This also extends to a functor: for a degree $0$ right $S$-module map $u:M\to N$, define $u^{\grstar}:N^{\grstar}\to M^{\grstar}$ by
\[
(u^{\grstar})_i:(N_{-i})^*\longrightarrow(M_{-i})^*,
\qquad f\longmapsto f\circ u|_{M_{-i}}
\quad(i\in\Z).
\]
For a degree $0$ left module map $u:M\to N$, the same formula defines ${}^{\grstar}u:{}^{\grstar}N\to{}^{\grstar}M$.
With the actions from Definition~\ref{def:S-duals}, we obtain functors
\begin{equation}\label{eq:graded-S-dual-functors}
(-)^{\grstar}:(\Grmod S)^\op\longrightarrow\Grmod(S^\op),\qquad 
{}^{\grstar}(-):(\Grmod(S^\op))^\op\longrightarrow\Grmod S.
\end{equation}

Each functor in \eqref{eq:graded-S-dual-functors} restricts to locally finite-dimensional modules: replace $\Grmod$ by $\grmod$ in its source and target.
Indeed, for $R=S$ or $S^\op$ and $M\in\grmod R$,
\begin{equation}\label{eq:graded-S-dual-finiteness}
\dim_\Bbbk\Hom_R(M_{-i},R)
\leq(\dim_\Bbbk M_{-i})(\dim_\Bbbk R)<\infty
\quad(i\in\Z).
\end{equation}

\begin{lemma}\label{lem:S-bidual}
We have dualities
\[
\begin{aligned}
(-)^*:(\mod S)^\op & \xrightarrow{\sim}\mod(S^\op), & {}^*(-):(\mod(S^\op))^\op &\xrightarrow{\sim}\mod S,\\
(-)^{\grstar}:(\grmod S)^\op&\xrightarrow{\sim}\grmod(S^\op), &
{}^{\grstar}(-):(\grmod(S^\op))^\op&\xrightarrow{\sim}\grmod S.
\end{aligned}
\]
\end{lemma}

\begin{proof}
We show that evaluation gives natural isomorphisms of functors
\[
\id_{\mod S}\xrightarrow{\sim}{}^*((-)^*),
\qquad
\id_{\mod(S^\op)}\xrightarrow{\sim}({}^*(-))^*,
\]
with components
\[
\begin{aligned}
\mathrm{ev}_M:M&\xrightarrow{\sim}{}^*(M^*),
& m&\longmapsto(f\mapsto f(m))
&& (M\in\mod S),\\
\mathrm{ev}_N:N&\xrightarrow{\sim}({}^*N)^*,
& n&\longmapsto(g\mapsto g(n))
&& (N\in\mod(S^\op)).
\end{aligned}
\]
For $u:M\to N$ in $\mod S$, $m\in M$, and $f\in N^*$,
\[
\bigl({}^*(u^*)(\mathrm{ev}_M(m))\bigr)(f)
=(u^*f)(m)=f(u(m))=\mathrm{ev}_N(u(m))(f).
\]
Hence $ {}^*(u^*)\circ\mathrm{ev}_M=\mathrm{ev}_N\circ u$.
For $M=S$, evaluation is an isomorphism: its inverse sends $F\in{}^*(S^*)$ to $F(\id_S)$.
Indeed, every $f\in S^*$ satisfies $f(s)=f(1)s$, so $F(f)=f(1)F(\id_S)=f(F(\id_S))$.
Restriction to summands identifies the dual of a finite direct sum with the direct sum of the duals; evaluation therefore is an isomorphism for $S^{\oplus n}$.
By Corollary~\ref{cor:semisimple-summands}, $M\oplus L\cong S^{\oplus n}$ for some $L$.
Under the direct sum identifications, $\mathrm{ev}_{M\oplus L}=\mathrm{ev}_M\oplus\mathrm{ev}_L$, so $\mathrm{ev}_M$ is an isomorphism.
Replacing $S$ by $S^\op$ proves the left module statement.

For the graded cases, apply the ungraded case degree-wise.  Note that since $S$ is finite-dimensional, finitely generated $S$-module is the same as finite-dimensional $S$-modules.
\end{proof}

For a bimodule ${}_SM_S$, the linear and one-sided duals retain both actions:
\begin{equation}\label{eq:bimodule-dual-actions}
(s\lambda t)(m)=\lambda(tms),\qquad
(sft)(m)=sf(tm),\qquad
(sgt)(m)=g(ms)t,
\end{equation}
where $s,t\in S$, $m\in M$, $\lambda\in D(M)$, $f\in\Hom_S(M,S)$, and $g\in\Hom_{S^\op}(M,S)$.
The functionals still take values in $\Bbbk$ or $S$.
With precomposition on morphisms, we obtain
\begin{equation}\label{eq:one-sided-bimodule-dual-functors}
\begin{aligned}
D,\ (-)^*,\ {}^*(-)&:(\Mod S^\en)^\op\longrightarrow\Mod S^\en,\\
\bbD,\ (-)^{\grstar},\ {}^{\grstar}(-)&:(\Grmod S^\en)^\op\longrightarrow\Grmod S^\en.
\end{aligned}
\end{equation}
The restrictions to finite-dimensional modules, and to locally finite-dimensional graded modules, are dualities by Lemmas~\ref{lem:linear-bidual} and~\ref{lem:S-bidual}.
Evaluation preserves both actions: for $m\in M$, $f\in M^*$, and $s,t\in S$,
\[
(s\,\mathrm{ev}_M(m)\,t)(f)
=f(sm)t=f(smt)=\mathrm{ev}_M(smt)(f).
\]
The corresponding checks for $D$ and the left dual use the same evaluation formula; applying them in each degree gives the graded statements.

Recall from Definition~\ref{def:enveloping-algebra} and \eqref{eq:enveloping-action} that
$S^\en=S\otimes_\Bbbk S^\op$ and $S$-bimodules are the same thing as $S^\en$-modules via $m\cdot(s\otimes t^\op)=tms$.

\begin{definition}\label{def:en-dual}
The \defn{two-sided dual} of a bimodule ${}_SM_S$ is
\[
M^{*_{\en}}:=\Hom_{S^{\en}}(M,S^{\en}),
\qquad (sFt)(m):=(s\otimes t^\op)F(m)
\quad(s,t\in S,\ F\in M^{*_{\en}},\ m\in M).
\]
\end{definition}

In Definition~\ref{def:en-dual}, $M$ has the right action \eqref{eq:enveloping-action}, so $F(tms)=F(m)(s\otimes t^\op)$.
The output is a left $S^{\en}$-module by multiplication on the values, which gives the displayed bimodule action.
The two-sided dual takes values in $S^{\en}$, not in $S$.
On bimodule morphisms it acts by precomposition, as do $D$, $(-)^*$, and ${}^*(-)$; each is a contravariant functor on the category of $S$-bimodules.
The two-sided dual records both module actions in the target algebra $S^\en$.

\phantomsection\label{not:graded-en-dual}
For a graded bimodule ${}_SM_S$, define its \defn{graded two-sided dual} by
\[
M^{\grstar[_{\en}]}:=\bigoplus_{i\in\Z}(M_{-i})^{*_{\en}},
\qquad (M^{\grstar[_{\en}]})_i=(M_{-i})^{*_{\en}}.
\]
Definition~\ref{def:en-dual} gives a functor
\begin{equation}\label{eq:graded-bimodule-dual-functors}
(-)^{\grstar[_{\en}]}:(\Grmod S^\en)^\op\longrightarrow\Grmod S^\en.
\end{equation}
On a degree $0$ bimodule map $u:M\to N$, its degree $i$ component is
\[
(u^{\grstar[_{\en}]})_i:(N_{-i})^{*_{\en}}\longrightarrow(M_{-i})^{*_{\en}},
\qquad F\longmapsto F\circ u|_{M_{-i}}
\quad(i\in\Z).
\]
The functor restricts to locally finite-dimensional bimodules by \eqref{eq:graded-S-dual-finiteness} with $R=S^\en$.

The four duals above follow the distinction in Grant--Iyama \cite[\S2.1, pp.~2592--2593]{GI20}, with the following notation changes:
\[
\begin{array}{c|c|c}
\text{construction}&\text{Grant--Iyama}&\text{these notes}\\ \hline
\Bbbk\text{-linear dual}&M^*&D(M)\\
\text{left }S\text{-linear dual}&M^{*_\ell}&{}^*M\\
\text{right }S\text{-linear dual}&M^{*_r}&M^*\\
\text{two-sided dual}&M^\vee&M^{*_{\en}}
\end{array}
\]
The last row also changes the convention for enveloping modules.
Their $F:M\to S^\en$ is left $S^\en$-linear.
Define
\[
\sigma:S^\en\longrightarrow S^\en,\qquad
a\otimes b^\op\longmapsto b\otimes a^\op.
\]
Since $\sigma(xy)=\sigma(y)\sigma(x)$ for $x,y\in S^\en$, the map $\sigma\circ F:M\to S^\en$ satisfies
\[
(\sigma\circ F)(tms)=(\sigma\circ F)(m)(s\otimes t^\op)
\quad(s,t\in S,\ m\in M),
\]
and belongs to our $M^{*_{\en}}$.
Thus $M^\vee\to M^{*_{\en}}$, $F\mapsto\sigma\circ F$, translates their two-sided dual into our convention.
In particular, their star denotes the linear dual; our star denotes the right $S$-dual.

\subsection{Comparing the duals}

The one-sided $S$-duals and the $\Bbbk$-linear dual are naturally isomorphic after choosing a symmetrising form on $S$.



\begin{definition}\label{def:symmetrising-form}
A \defn{symmetrising form} on a finite-dimensional $\Bbbk$-algebra $S$ is a linear map $\tau:S\to\Bbbk$ satisfying
\[
\tau(st)=\tau(ts)\quad\text{for all }s,t\in S,
\qquad
\bigl(\tau(st)=0\text{ for all }t\in S\bigr)\ \Longrightarrow\ s=0.
\]
An algebra admitting such a form is \defn{symmetric}.
\end{definition}

The second condition in Definition~\ref{def:symmetrising-form} says that the bilinear form $(s,t)\mapsto\tau(st)$ is \defn{non-degenerate}: every non-zero $s\in S$ pairs non-trivially with some $t\in S$.

Every finite-dimensional semisimple $\Bbbk$-algebra admits a symmetrising form.
For products of matrix algebras over $\Bbbk$, the forms are explicit:

\begin{example}\label{ex:symmetrising-form}
The following maps are symmetrising forms:
\[
\begin{aligned}
\tau:\Bbbk^{\oplus n}&\longrightarrow\Bbbk,
& (a_1,\ldots,a_n)&\longmapsto\sum_{j=1}^n a_j,\\
\tau:\Mat_n(\Bbbk)&\longrightarrow\Bbbk,
& a&\longmapsto\Tr(a):=\sum_{i=1}^n a_{ii},\\
\text{and more generally,}\quad \tau:\prod_{j=1}^n\Mat_{n_j}(\Bbbk)&\longrightarrow\Bbbk,
& (a_1,\ldots,a_n)&\longmapsto\sum_{j=1}^n\Tr(a_j).
\end{aligned}
\]
\end{example}

Generally, any finite-dimensional semisimple algebra has a symmetrising form, by Eilenberg--Nakayama.\footnote{S.~Eilenberg and T.~Nakayama, \emph{On the dimension of modules and algebras, II. Frobenius algebras and quasi-Frobenius rings}, Nagoya Math. J. \textbf{9} (1955), 1--16, \S2, Propositions 3--5, pp.~4--5; \href{https://doi.org/10.1017/S0027763000023229}{doi:10.1017/S0027763000023229}.}
For the graded analogue, see \cite{DNN17}.
For example, let $\lambda_j:\Delta_j\to\Bbbk$ be a symmetrising form on a finite-dimensional division $\Bbbk$-algebra $\Delta_j$ for $j=1,\ldots,n$.
Then the semisimple algebra $S=\prod_{j=1}^n\Mat_{n_j}(\Delta_j)$ has a symmetrising form given by
\[
\tau:S\longrightarrow\Bbbk,\qquad
(a_1,\ldots,a_n)\longmapsto
\sum_{j=1}^n\sum_{i=1}^{n_j}\lambda_j((a_j)_{ii}).
\]

\begin{theorem}[Comparison of duals]\label{prop:dual-comparison}\label{cor:graded-dual-comparison}
Choose a symmetrising form $\tau:S\to\Bbbk$.
For $M\in\mod S$ and $N\in\mod(S^\op)$, postcomposition with $\tau$ gives $S$-module isomorphisms
\begin{equation}\label{eq:dual-comparison}
\alpha_M:M^*\longrightarrow D(M),\quad f\longmapsto\tau\circ f,
\qquad
\beta_N:{}^*N\longrightarrow D(N),\quad g\longmapsto\tau\circ g,
\end{equation}
natural in $M$ and $N$.
Applied in each degree, the same maps give natural isomorphisms $M^{\grstar}\cong\bbD(M)$ and ${}^{\grstar}N\cong\bbD(N)$ for $M\in\grmod S$ and $N\in\grmod(S^\op)$; we again write $\alpha_M$ and $\beta_N$.
\end{theorem}

We can spell this out more precisely:


\underline{\emph{Ungraded duals.}}
The isomorphisms \eqref{eq:dual-comparison} give natural isomorphisms
\begin{equation}\label{eq:dual-comparison-functors}
\begin{aligned}
\alpha:(-)^*&\xrightarrow{\sim}D
&&\text{of functors }(\mod S)^\op\longrightarrow\mod(S^\op),\\
\beta:{}^*(-)&\xrightarrow{\sim}D
&&\text{of functors }(\mod(S^\op))^\op\longrightarrow\mod S.
\end{aligned}
\end{equation}

\underline{\emph{Graded duals.}}
Composition with $\tau$ in each degree gives natural isomorphisms
\begin{equation}\label{eq:graded-dual-comparison}
\begin{aligned}
\alpha:(-)^{\grstar}&\xrightarrow{\sim}\bbD
&&\text{of functors }(\grmod S)^\op\longrightarrow\grmod(S^\op),\\
\beta:{}^{\grstar}(-)&\xrightarrow{\sim}\bbD
&&\text{of functors }(\grmod(S^\op))^\op\longrightarrow\grmod S.
\end{aligned}
\end{equation}
For $M\in\grmod S$ and $N\in\grmod(S^\op)$, the degree $i$ components are $\alpha_{M_{-i}}$ and $\beta_{N_{-i}}$, respectively.

\begin{proof}
We prove the right module statement, i.e. taking $\alpha_M(f):= \tau\circ f$ yields 
\[
\text{ungraded case } \alpha: (-)^*\xrightarrow{\sim} D, \quad \text{and}\quad \text{graded case }\alpha: (-)^{(*)}\xrightarrow{\sim} \bbD.
\]
First, non-degeneracy of $\tau$ gives an isomorphism $S\to D(S)$, $a\mapsto(s\mapsto\tau(as))$.
Hence, for $\lambda\in D(M)$, there is a unique linear map $f:M\to S$ satisfying
\begin{equation}\label{eq:dual-inverses}
\tau(f(m)s)=\lambda(ms)\quad\text{for all }m\in M,\ s\in S.
\end{equation}
For $m\in M$ and $s,t\in S$,
\[
\tau(f(mt)s)=\lambda(mts)=\tau(f(m)ts),
\]
so non-degeneracy implies $f(mt)=f(m)t$.
Taking $s=1$ in \eqref{eq:dual-inverses} gives $\tau\circ f=\lambda$; conversely, every right $S$-linear preimage of $\lambda$ must satisfy \eqref{eq:dual-inverses}.
Thus $\alpha_M$ is bijective.

Both $M^*$ and $D(M)$ are left $S$-modules, so the component $\alpha_M:M^*\to D(M)$ must preserve the left action:
\[
\alpha_M(sf)(m)=\tau(sf(m))
=\tau(f(m)s)=\tau(f(ms))
=(s\alpha_M(f))(m)
\quad(s\in S,\ f\in M^*,\ m\in M).
\]
For a right module map $u:M\to N$ and $f\in N^*$, associativity of composition gives
$\alpha_M(f\circ u)=\alpha_N(f)\circ u$, proving naturality.

For the graded comparisons, apply the ungraded natural isomorphism degree-wise. 

Same argument applies dually to $S^\op$ and $\beta$.
\end{proof}

\begin{remark}\label{rem:standard-form}
The natural isomorphisms in Theorem~\ref{prop:dual-comparison} depend on the chosen symmetrising form.
For example, every choice $c_1,\ldots,c_n\in\Bbbk\setminus\{0\}$ gives a form
$\tau(a_1,\ldots,a_n)=\sum_{i=1}^n c_i a_i$ on $\Bbbk^{\oplus n}$.
Under Assumption~\ref{ass:basic-semisimple} we always use the \defn{standard symmetrising form} $\tau_S(\sum_{i=1}^na_ie_i)=\sum_{i=1}^na_i$, which involves no choice; then $\alpha_M$ and $\beta_N$ are canonical isomorphisms.
\end{remark}

\begin{remark}[General $S$]\label{rem:general-S-comparison}
For an arbitrary finite-dimensional semisimple $S$, identifying an $S$-dual with a $\Bbbk$-linear dual requires choosing a symmetrising form.
If $\tau'$ is another one, then $\tau'(s)=\tau(zs)$ for a central unit $z\in S$, and the two identifications differ by multiplication by $z$; see Exercise~\ref{ex:transported-dual-action}(3).
\end{remark}

\begin{remark}\label{rem:transported-dual-action}
Let $A$ be a graded algebra with $A_0=S$.
The space $\Hom_S(M,S)$ has no natural left $A$-module structure on its own, but acquires one after \emph{picking a symmetrising form $\tau$ of $S$}; under Assumption~\ref{ass:basic-semisimple} we pick $\tau=\tau_S$.
Indeed, the graded comparison $\alpha_M$ of Theorem~\ref{prop:dual-comparison} transports the canonical left $A$-action on $\bbD(M)$ to the graded $S$-dual of a locally finite-dimensional module:
\[
af=\alpha_M^{-1}(a\alpha_M(f))\quad\text{for right }M,
\qquad ga=\beta_N^{-1}(\beta_N(g)a)\quad\text{for left }N,
\]
where $a\in A$, $f\in M^{\grstar}$, and $g\in{}^{\grstar}N$.
Similarly for right action on left dual.
%For comparison, a bimodule ${}_AX_S$ gives a bimodule ${}_S\bigl(\Hom_S(X,S)\bigr)_A$ by$(sfa)(x)=s f(ax)$.
%This is the construction used in \cite[\S2.2]{BGS96}; it uses commuting actions on opposite sides, not merely restriction of a right $A$-module to $S$.
%Exercise~\ref{ex:transported-dual-action} compares two transported actions on one graded $S$-dual.
\end{remark}


The comparison theorem has a bimodule version:

\begin{theorem}
With the actions in \eqref{eq:bimodule-dual-actions}, we obtain natural isomorphisms of functors
\begin{equation}\label{eq:bimodule-dual-comparison}
\begin{aligned}
(-)^*\xrightarrow[\alpha]{\sim}D\xleftarrow[\beta]{\sim}{}^*(-)
&\quad:(\mod S^\en)^\op\longrightarrow\mod S^\en,\\
(-)^{\grstar}\xrightarrow[\alpha]{\sim}\bbD\xleftarrow[\beta]{\sim}{}^{\grstar}(-)
&\quad:(\grmod S^\en)^\op\longrightarrow\grmod S^\en.
\end{aligned}
\end{equation}

There is also a natural isomorphism $\omega:(-)^{*_\en}\xrightarrow{\sim}(-)^*$ given by
\begin{equation}\label{eq:en-dual-comparison}
\omega_M:M^{*_{\en}}\to M^*, \qquad F\mapsto \big(m\mapsto \sum_j a_j\tau(b_j)\big)
\quad\text{for}\quad
F(m)=\sum_j a_j\otimes b_j^\op,
\end{equation}
which also depends on $\tau$, but does not require $S^\en$ to be semisimple.

Consequently, on finite-dimensional $S$-bimodules and lfd graded $S$-bimodules we have natural isomorphisms:
\begin{align*}
\text{Ungraded:}\quad &(-)^{*_{\en}}\xrightarrow[\omega]{\sim}(-)^*
\xrightarrow[\alpha]{\sim}D\xleftarrow[\beta]{\sim}{}^*(-) \text{ as functors } (\mod S^\en)^\op\to\mod S^\en, \\
\text{Graded:}\quad  &(-)^{\grstar[_{\en}]}\xrightarrow[\omega]{\sim}(-)^{\grstar}
\xrightarrow[\alpha]{\sim}\bbD\xleftarrow[\beta]{\sim}{}^{\grstar}(-) \text{ as functors } (\grmod S^\en)^\op\to\grmod S^\en.
\end{align*}
\end{theorem}
\begin{proof}
The form $\tau_{\en}(a\otimes b^\op)=\tau(a)\tau(b)$ is symmetrising on $S^{\en}$.
Indeed, $\tau_{\en}((a\otimes b^\op)(c\otimes d^\op))=\tau(ac)\tau(db)=\tau(ca)\tau(bd)$, and its pairing is the tensor product of two non-degenerate pairings.
The formula \eqref{eq:en-dual-comparison} satisfies
\[
\omega_M(F)(ms)=\omega_M(F)(m)s,\qquad
\omega_M(sFt)=s\omega_M(F)t,\qquad
\alpha_M(\omega_M(F))=\tau_{\en}\circ F.
\]
The argument \eqref{eq:dual-inverses}, applied to $S^{\en}$ and $\tau_{\en}$, makes $F\mapsto\tau_{\en}\circ F$ an isomorphism.
Thus $\omega_M$ is an isomorphism.
For a bimodule map $u:M\to N$ and $F\in N^{*_{\en}}$, the formula \eqref{eq:en-dual-comparison} gives $\omega_M(F\circ u)=\omega_N(F)\circ u$, proving naturality.
Taking the direct sum of $\omega_{M_{-i}}$ over $i\in\Z$ gives the graded two-sided comparison.
Each component is a bimodule isomorphism, and naturality follows in each degree from the ungraded comparison.
\end{proof}
\end{advanced}


\subsection{Exercises}

\begin{exercise}[Matrix calculations]
Let $n\geq2$ and $A=\Bbbk A_n$.
Verify the grading in Example~\ref{ex:triangular-grading} using $e_{ij}e_{pq}=\delta_{jp}e_{iq}$.
Compute $\rad A$ from its definition as the intersection of all maximal right ideals.

\emph{Hint.} The strictly upper triangular ideal $J$ satisfies $J^n=0$.
Show that $LJ=0$ for every simple right $A$-module $L$, and use the diagonal maps $A\to\Bbbk$, $X\mapsto X_{ii}$, for $1\leq i\leq n$.
\end{exercise}

\begin{exercise}[Shifts and duals]
Let $A=\Bbbk A_3$ and $M=e_{22}A$ with the inherited grading.
\begin{enumerate}
\item Find a homogeneous basis of $\bbD(M)$, its degrees, and all non-zero actions of matrix units on the basis vectors.
\item Show that $\Bbbk e_{13}\subset Ae_{33}$ is a graded left submodule and that
\[
\bbD(M)\cong(Ae_{33}/\Bbbk e_{13})(1)
\]
as graded left $A$-modules.
\item Deduce that $\bbD(M)\not\cong(Ae_{jj})(r)$ for every $j\in\{1,2,3\}$ and $r\in\Z$.
\end{enumerate}
\emph{Hint for (3).} Compare dimensions and the actions of $e_{11},e_{22},e_{33}$.
\end{exercise}

\begin{exercise}[Homogeneous maps]\label{ex:infinite-hom}
The finite-generation hypothesis in Lemma~\ref{lem:homogeneous-hom-sum} is needed.
Let $S=\Bbbk e_1\oplus\Bbbk e_2$ with componentwise multiplication, let $N=e_1S$ have degree $0$, and put
\[
M=\bigoplus_{i\geq0}(e_1S)(-i)=\bigoplus_{i\geq0}\Bbbk m_i,
\qquad \deg m_i=i,\quad m_i e_1=m_i,\quad m_i e_2=0.
\]
For the $S$-linear map $f:M\to N$ with $f(m_i)=e_1$ for every $i\geq0$, determine all homogeneous components and prove that $f$ is not a finite sum of homogeneous maps.
\end{exercise}

\begin{hidden}
\begin{proof}[Solution]
For a finitely supported family $(c_i)_{i\geq0}$,
\[
f\Bigl(\sum_{i\geq0}c_i m_i\Bigr)=\sum_{c_i\neq0}c_i e_1.
\]
On the summand $(e_1S)(-i)$, the restriction $f|_{(e_1S)(-i)}:(e_1S)(-i)\to e_1S$ is the underlying identity map, of degree $-i$.
Explicitly, the homogeneous components are
\[
f_{-r}:M\longrightarrow N,\qquad f_{-r}(m_i)=\delta_{ri}e_1
\quad\text{for }r,i\geq0,
\]
and $f_d=0$ for $d>0$.
A homogeneous map $M\to N$ of degree $d$ vanishes on $m_i$ unless $d=-i$.
Every finite sum of homogeneous maps therefore vanishes on all but finitely many $m_i$, whereas $f(m_i)=e_1$ for every $i\geq0$.
Hence $f$ is not a finite sum of homogeneous maps.
\end{proof}
\end{hidden}

\begin{advanced}
\begin{exercise}[Dependence on the form]\label{ex:dual-blocks}
Let $S=(\Bbbk A_2)_0=\Bbbk e_{11}\oplus\Bbbk e_{22}$, $V=(\Bbbk A_2)_1=\Bbbk e_{12}$, and $c_1,c_2\in\Bbbk\setminus\{0\}$.
Set $\tau(ae_{11}+be_{22})=c_1a+c_2b$ and let $\varepsilon:V\to\Bbbk$ satisfy $\varepsilon(e_{12})=1$.
Find the right and left $S$-linear maps corresponding to $\varepsilon$ under \eqref{eq:dual-comparison}, and determine how they depend on $c_1,c_2$.
\end{exercise}

\begin{hidden}
\begin{proof}[Solution]
Right $S$-linearity forces $f(e_{12})=f(e_{12})e_{22}\in\Bbbk e_{22}$, and left $S$-linearity forces $g(e_{12})=e_{11}g(e_{12})\in\Bbbk e_{11}$.
Thus $\tau\circ f=\varepsilon=\tau\circ g$ requires
\[
f(e_{12})=c_2^{-1}e_{22},\qquad g(e_{12})=c_1^{-1}e_{11}.
\]
Conversely, the linear maps with these values satisfy, for $s=ae_{11}+be_{22}\in S$,
\[
f(e_{12}s)=bc_2^{-1}e_{22}=f(e_{12})s,\qquad
g(se_{12})=ac_1^{-1}e_{11}=sg(e_{12}).
\]
Their compositions with $\tau$ take $e_{12}$ to $1$, as required.
The original worked example is the case $c_1=c_2=1$.
\end{proof}
\end{hidden}

\begin{exercise}[Dependence of the action on the form]\label{ex:transported-dual-action}
The transported $A$-action can depend on the symmetrising form, but its graded isomorphism class does not.
Let $A=\Bbbk A_2$, $S=\Bbbk e_{11}\oplus\Bbbk e_{22}$, and $M=e_{11}A$, with the inherited gradings.
For $c=(c_1,c_2)\in(\Bbbk\setminus\{0\})^2$, define
\begin{equation}\label{eq:transported-dual-form}
\tau_c:S\longrightarrow\Bbbk,\qquad
a e_{11}+b e_{22}\longmapsto c_1a+c_2b
\quad\text{for all }a,b\in\Bbbk.
\end{equation}
Write $\alpha_c:M^{\grstar}\to\bbD(M)$, $f\mapsto\tau_c\circ f$, and let $a\cdot_c f=\alpha_c^{-1}(a\alpha_c(f))$ be the transported action of Remark~\ref{rem:transported-dual-action}.
Consider the right $S$-linear maps
\begin{equation}\label{eq:transported-dual-functionals}
\begin{aligned}
\varphi_0:M&\longrightarrow S,& a e_{11}+b e_{12}&\longmapsto a e_{11},\\
\varphi_1:M&\longrightarrow S,& a e_{11}+b e_{12}&\longmapsto b e_{22}
\quad\text{for all }a,b\in\Bbbk.
\end{aligned}
\end{equation}
\begin{enumerate}
\item Verify that $\tau_c$ is symmetrising and that $\varphi_0,\varphi_1$ form a homogeneous basis of $M^{\grstar}$. Determine their degrees.
\item Calculate the actions of $e_{11},e_{22},e_{12}$ on $\varphi_0,\varphi_1$.
For $c,c'\in(\Bbbk\setminus\{0\})^2$, prove that $\cdot_c=\cdot_{c'}$ if and only if $c'=\rho c$ for some $\rho\in\Bbbk\setminus\{0\}$.
\item More generally, let $A$ be a locally finite-dimensional graded algebra with semisimple $S=A_0$, and let $M$ be a locally finite-dimensional graded right $A$-module.
Let $\tau,\tau'$ be symmetrising forms on $S$, with comparison maps $\alpha_\tau,\alpha_{\tau'}:M^{\grstar}\to\bbD(M)$ given by composition with $\tau,\tau'$, respectively, and transported actions $\cdot_\tau,\cdot_{\tau'}$.
Prove that there is an invertible $z\in S$ satisfying $\tau'(s)=\tau(zs)$ for every $s\in S$, and that
\[
\gamma:(M^{\grstar},\cdot_{\tau'})\longrightarrow(M^{\grstar},\cdot_\tau),
\qquad (\gamma f)(m)=z f(m)
\]
is an isomorphism of graded left $A$-modules.
Deduce that the modules in (2) are isomorphic for every $c,c'$, and give the isomorphism on $\varphi_0,\varphi_1$.
\end{enumerate}
\emph{Hint for (3).} Compare $\alpha_{\tau'}$ with $\alpha_\tau$.
\end{exercise}

\begin{hidden}
\begin{proof}[Solution]
\emph{The form and the basis.}
The algebra $S$ is commutative, so $\tau_c(st)=\tau_c(ts)$.
If $s=ae_{11}+be_{22}$ satisfies $\tau_c(se_{11})=\tau_c(se_{22})=0$, then $c_1a=c_2b=0$, hence $s=0$.
Thus $\tau_c$ is non-degenerate.

A right $S$-linear map $f:M\to S$ must satisfy $f(e_{11})\in Se_{11}=\Bbbk e_{11}$ and $f(e_{12})\in Se_{22}=\Bbbk e_{22}$.
Conversely, any choices of these two values define such a map.
Consequently, $\varphi_0,\varphi_1$ form a basis, with degrees $0,-1$, respectively.

\emph{The actions.}
For $m=ae_{11}+be_{12}$,
\[
\alpha_c(\varphi_0)(m)=c_1a,\qquad
\alpha_c(\varphi_1)(m)=c_2b,\qquad
\alpha_c(e_{12}\cdot_c\varphi_1)(m)
=\alpha_c(\varphi_1)(me_{12})=c_2a.
\]
Using also $me_{11}=ae_{11}$ and $me_{22}=be_{12}$ gives all non-zero actions:
\[
e_{11}\cdot_c\varphi_0=\varphi_0,\qquad
e_{22}\cdot_c\varphi_1=\varphi_1,\qquad
e_{12}\cdot_c\varphi_1=\frac{c_2}{c_1}\varphi_0.
\]
All other actions of the three matrix units on the two basis vectors are zero.
Consequently, $\cdot_c=\cdot_{c'}$ if and only if $c_2/c_1=c'_2/c'_1$.
This equality is equivalent to $c'=\rho c$ with $\rho=c'_1/c_1\neq0$.

\emph{The general isomorphism.}
Non-degeneracy of $\tau$ makes
\[
S\longrightarrow D(S),\qquad z\longmapsto\bigl(s\mapsto\tau(zs)\bigr)
\]
injective, hence bijective since $S$ is finite-dimensional.
Choose $z$ mapping to $\tau'$.
If $zs=0$, then $\tau'(st)=\tau(zst)=0$ for every $t\in S$, so $s=0$ by non-degeneracy of $\tau'$.
Thus multiplication by $z$ on the left is bijective.
Choose $u\in S$ with $zu=1$; injectivity applied to $z(uz-1)=0$ gives $uz=1$.

For $f\in M^{\grstar}$, $m\in M$, and $s\in S$, one has $zf(ms)=zf(m)s$, so $\gamma f$ is right $S$-linear.
Multiplication by $z$ preserves degrees, and multiplication by $z^{-1}$ gives the inverse of $\gamma$.
Moreover, $\alpha_\tau\gamma=\alpha_{\tau'}$, so for every $a\in A$ and $f\in M^{\grstar}$,
\[
\gamma(a\cdot_{\tau'}f)
=\alpha_\tau^{-1}\bigl(a\alpha_{\tau'}(f)\bigr)
=a\cdot_\tau\gamma(f).
\]
Hence $\gamma$ is a graded left $A$-module isomorphism.
For the forms $\tau_c,\tau_{c'}$, take
\[
z=\frac{c'_1}{c_1}e_{11}+\frac{c'_2}{c_2}e_{22}.
\]
The resulting isomorphism $(M^{\grstar},\cdot_{c'})\to(M^{\grstar},\cdot_c)$ sends
$\varphi_0\mapsto(c'_1/c_1)\varphi_0$ and $\varphi_1\mapsto(c'_2/c_2)\varphi_1$.
\end{proof}
\end{hidden}
\end{advanced}

\pagebreak

